% !TeX root = ../main.tex

\chapter{Discussion}\label{chapter:discussion}

% ~2--3 pages

\section{Trade-Off Summary}\label{sec:tradeoff-summary}
\begin{itemize}
  \item Restate the RQ and summarize the answer with concrete numbers from Chapter 5
  \item Optimal allocation of CI budget across test layers
  \item The FixedEmbedder as a key enabler: percentage of integration tests that can run GPU-free; what fraction of real bugs does it catch vs.\ miss?
  \item Diminishing returns: at what point does adding more metamorphic tests stop improving fault detection?
  \item Cost of false security: what types of bugs can only be caught with real GPU inference?
\end{itemize}

\section{Generalizability}\label{sec:generalizability}
\begin{itemize}
  \item Beyond biocentral: applicability to other PLM-based services (e.g., ColabFold, bio-embeddings servers)
  \item Beyond protein embeddings: applicability to NLP embedding services (sentence transformers, BERT-based APIs)
  \item Beyond embeddings: applicability to other ML inference services (image classifiers, recommendation systems)
  \item Which parts of the approach are domain-specific (canonical protein dataset, AA mutation MRs) vs.\ generalizable (test pyramid, FixedEmbedder pattern, oracle framework)?
  \item Multi-service testing: the approach handles complex Docker deployments (server + workers + GPU + DB) --- relevant for any microservice-based ML system
\end{itemize}

\section{Limitations}\label{sec:limitations}
\begin{itemize}
  \item Single PLM for experiments: ESM2-T6-8M is the smallest ESM2 variant; larger models (650M, 3B) or different architectures (ProtT5) may exhibit different behavior
  \item No adversarial testing: the metamorphic experiments use structured transformations, not adversarial inputs designed to fool models
  \item FixedEmbedder expressiveness: SHA-256-based embeddings do not capture real biological properties --- bugs that require biologically meaningful embeddings will be missed
  \item No end-user study: the test suite was developed by the author; no external developer usability evaluation
  \item Limited mutation testing scope: mutmut only targets unit tests; integration/property tests not evaluated for mutation score
  \item Hardware-specific timings: execution times depend on specific CPU/GPU; relative ratios are more meaningful than absolutes
\end{itemize}

\section{Threats to Validity}\label{sec:threats}
\begin{itemize}
  \item Internal validity:
    \begin{itemize}
      \item GPU non-determinism tolerance thresholds were chosen empirically --- different thresholds could change pass/fail outcomes
      \item FixedEmbedder monkey-patching may mask real HTTP/serialization bugs
      \item Mutation testing results depend on mutmut's mutation operators --- different tools may produce different scores
    \end{itemize}
  \item External validity:
    \begin{itemize}
      \item Results specific to biocentral's architecture (FastAPI + Redis + Triton) --- other frameworks may have different failure modes
      \item ESM2-T6-8M behavior may not generalize to larger or differently trained models
      \item UniRef50 sequences are biased toward well-studied proteins --- rare/novel proteins underrepresented
    \end{itemize}
  \item Construct validity:
    \begin{itemize}
      \item ``Test coverage'' measured via mutation score and code coverage --- other metrics (fault detection rate in production) would be more meaningful but are unavailable
      \item ``CI execution time'' measured on a single machine --- cloud CI environments may differ
    \end{itemize}
\end{itemize}
